% PDF semantic block environments — academic monograph style
% Matched to the HTML CSS in styles/custom.css (Section 12).

\usepackage{xcolor}
\usepackage{tcolorbox}
\tcbuselibrary{skins,breakable}
\usepackage{tocvsec2}
\usepackage{multicol}

% Suppress indent of the paragraph immediately following a semantic box.
% \noindent\ignorespaces glued directly onto \end{semanticX} — safe against
% both mdframed (this file's first box implementation) and tcolorbox (the
% current one, switched to below for its native per-side rule-width
% support), tested in isolated xelatex reproductions against each.
\newcommand{\suppressnextindent}{\noindent\ignorespaces}

% Footnotes flush left, no indent — on request. LaTeX's kernel default
% (\@makefntext) right-aligns the footnote mark inside an invisible
% 1.8em box (\hb@xt@1.8em{\hss\@makefnmark}), which reads as a small
% hanging indent before the mark and text both start. This drops that
% box entirely: mark then text, flush against the true left margin.
\makeatletter
\renewcommand{\@makefntext}[1]{%
  \setlength{\parindent}{0pt}%
  \noindent\@makefnmark\ #1%
}
\makeatother

% Per-type border colors — the exact light-mode --book-sb-*-border hex
% values from styles/themes/procyon.css, so each box reads as the same
% color family as its HTML counterpart. premise/argument/untyped share
% definition's own gray on the website too (their own --book-sb-*-border
% values are set to `var(--book-sb-definition-border)`), so they do here
% as well — not separate colors by omission, matching that inheritance
% intentionally. Also reused directly for the native "Note" callout
% restyle below (--book-sb-untyped-border is the same #6B6B6B gray).
\definecolor{sbDefinition}{HTML}{6B6B6B}
\definecolor{sbCondition}{HTML}{2A5F9C}
\definecolor{sbIdentification}{HTML}{3A6432}
\definecolor{sbPrinciple}{HTML}{B4552E}
\definecolor{sbLemma}{HTML}{7A2C9E}
\definecolor{sbCorollary}{HTML}{1E6E78}
\definecolor{sbTheorem}{HTML}{B8860B}
\definecolor{sbQuickref}{HTML}{6B6560}

% Shared box chassis — copied directly from Quarto's OWN native callout
% geometry (Tip/Warning/Important/Caution, unchanged, immediately below
% in this file), not approximated: on request, after several passes that
% kept missing the mark ("don't they just have a style you can almost
% copy... the only thing I want different is literally the colors and
% the no icon") — read the ACTUAL parameters Quarto generates for its own
% callouts straight out of this book's own compiled .tex (e.g. for Note:
% "enhanced jigsaw, arc=.35mm, bottomrule=.15mm, bottomtitle=1mm,
% colback=white, colbacktitle=COLOR!10!white, colframe=COLOR-frame,
% coltitle=black, left=2mm, leftrule=.75mm, rightrule=.15mm,
% title=ICON+TEXT, titlerule=0mm, toprule=.15mm, toptitle=1mm") and
% matched every one of those values here — every earlier pass had
% guessed at reasonable-looking numbers instead (9pt/7pt/6pt padding, a
% separate borderline-west overlay for the left accent) and NEVER
% actually matched Quarto's own numbers, which is exactly why the result
% kept reading as "still not there" no matter how many individual
% properties got nudged. leftrule alone (thicker than the other three
% sides) is what gives the left accent — no separate overlay decoration
% needed, and titlerule=0mm (Quarto's own value: no separate rule
% distinct from the frame's own bottomrule closing the title band).
% Only real, deliberate departures from Quarto's own values: title=#1
% (dynamic per-instance text/argument, not a fixed icon+label), and
% colbacktitle at 15% instead of Quarto's 10% (this book's own per-type
% colors are less saturated than Quarto's Bootstrap-derived theme colors,
% so 10% read as barely-there in an earlier check — 15% is the smallest
% bump that stayed visible).
\tcbset{
  semanticboxbase/.style={
    enhanced,
    arc=0.35mm,
    toprule=0.15mm, bottomrule=0.15mm, rightrule=0.15mm, leftrule=0.75mm,
    colback=white,
    coltitle=black,
    left=2mm, right=2mm, top=2mm, bottom=2mm,
    toptitle=1mm, bottomtitle=1mm,
    titlerule=0mm,
    before skip=10pt, after skip=10pt,
    before upper={\setlength{\parindent}{0pt}},
  }
}

\newtcolorbox{semanticboxdefinition}[1]{semanticboxbase, colframe=sbDefinition, colbacktitle=sbDefinition!15!white, title=#1}
\newtcolorbox{semanticboxtheorem}[1]{semanticboxbase, colframe=sbTheorem, colbacktitle=sbTheorem!15!white, title=#1}
\newtcolorbox{semanticboxlemma}[1]{semanticboxbase, colframe=sbLemma, colbacktitle=sbLemma!15!white, title=#1}
\newtcolorbox{semanticboxcorollary}[1]{semanticboxbase, colframe=sbCorollary, colbacktitle=sbCorollary!15!white, title=#1}
\newtcolorbox{semanticboxcondition}[1]{semanticboxbase, colframe=sbCondition, colbacktitle=sbCondition!15!white, title=#1}
\newtcolorbox{semanticboxquickreference}[1]{semanticboxbase, colframe=sbQuickref, colbacktitle=sbQuickref!15!white, title=#1}
\newtcolorbox{semanticboxidentification}[1]{semanticboxbase, colframe=sbIdentification, colbacktitle=sbIdentification!15!white, title=#1}
\newtcolorbox{semanticboxprinciple}[1]{semanticboxbase, colframe=sbPrinciple, colbacktitle=sbPrinciple!15!white, title=#1}

% Hanging-indent bibliography — PDF equivalent of the HTML
% body.references-page CSS rule (styles/custom.css, section 15). Applied by
% filters/semantic-blocks.lua to the chapter titled "References". First line
% of each entry sits flush left; wrapped lines indent 1.5em, instead of
% LaTeX's default first-line indent. \raggedright matches the HTML rule's
% text-align: left — ragged right edge, not fully justified.
\newenvironment{hangingrefs}{%
  \begingroup
  \setlength{\parindent}{0pt}%
  \setlength{\parskip}{4pt}%
  \everypar{\hangindent=1.5em\hangafter=1\relax}%
  \raggedright
}{%
  \par\endgroup
}

% Zero first-line indent on every paragraph — PDF equivalent of the HTML
% body.glossary-page CSS rule (styles/custom.css, section 16). Applied by
% filters/semantic-blocks.lua to the Glossary chapter, a flat list of
% one-paragraph entries where neither the book's normal first-line indent
% nor a hanging indent belongs — every entry sits flush left, ragged right
% (\raggedright), not fully justified, matching the HTML rule.
\newenvironment{flushparagraphs}{%
  \begingroup
  \setlength{\parindent}{0pt}%
  \setlength{\parskip}{5pt}%
  \raggedright
}{%
  \par\endgroup
}

% Two-column, smaller-text subject index — PDF equivalent of the HTML
% body.subject-index-page CSS rule (styles/custom.css, section 16b), with
% the standard back-of-book-index layout (two columns, reduced size) added
% on top. Applied by filters/semantic-blocks.lua to the Index chapter. Each
% entry's first line sits flush left; a wrapped continuation line nests in
% by 1em, matching the same 1em hanging indent used on the HTML index page.
% \raggedright matches the HTML rule's text-align: left.
\newenvironment{twocolindex}{%
  \begingroup
  \small
  \setlength{\parindent}{0pt}%
  \everypar{\hangindent=1em\hangafter=1\relax}%
  \raggedright
  \begin{multicols}{2}
}{%
  \end{multicols}
  \endgroup
}

% Native Quarto "Note" callout — the only callout type actually used in
% the book (the other four are kept in style-preview.qmd purely for
% reference) — restyled in filters/semantic-blocks.lua to use the same
% chassis as every semantic box above (Definition/Untyped's shared gray:
% title band, left accent, full-width rule, white body), instead of
% Quarto's own tcolorbox jigsaw/icon/title-band style. On request: "the
% note too... it needs to have the same kind of thing." Callouts can run
% long in normal prose, unlike the short structured semantic types above,
% so — unlike semanticboxbase — this one stays breakable.
\newtcolorbox{semanticboxnote}[1]{semanticboxbase, breakable, colframe=sbDefinition, colbacktitle=sbDefinition!15!white, title=#1}

% Tables — most tables (simple shape: one header row, one body, no
% multi-paragraph cells, no pandoc-assigned fractional column widths) are
% now hand-emitted directly by filters/semantic-blocks.lua's Table()
% function as a bordered longtable matching the same chassis as the
% semantic boxes and Note: its own existing header row tinted gray via
% colortbl's \rowcolor (not a separate title row), outer-edge vertical
% bars, \hline rules — see that function's own header comment for why
% this couldn't just be a tcolorbox wrapped around the existing output
% (longtable's cross-page breaking doesn't survive being nested in any
% other box). Any table that function declines (complex shape) falls
% through to this booktabs-based styling instead, applying to Pandoc's
% own default \begin{longtable} output: \heavyrulewidth/\lightrulewidth
% thin the default booktabs rules, and every longtable (both kinds) gets
% shrunk to \small, closer to the website's smaller table font-size
% (0.83rem vs 1rem body).
\usepackage{booktabs}
\usepackage{colortbl}
\usepackage{etoolbox}
\setlength{\heavyrulewidth}{0.9pt}
\setlength{\lightrulewidth}{0.5pt}
\AtBeginDocument{\AtBeginEnvironment{longtable}{\small}}
